%% ============================================================
%%  Exercises part -- opener for homework/problem families.
%%  Subfile of main.tex.
%% ============================================================
\documentclass[main.tex]{subfiles}
\begin{document}

\beginunittoc{Problem Set Conventions}{PROBLEM SETS}{}{Problem Set Conventions}{Problem Set Conventions}

\begin{partcontents}[7.5em]{Problem Sets in This Part}
  \partcontentitem{hw-1}{Homework 1}{The Real Numbers}
  \partcontentitem{hw-2}{Homework 2}{The Topology of Metric Spaces}
  \partcontentitem{hw-3}{Homework 3}{Sequences, Series, and Continuity}
\end{partcontents}

\vspace{0.65\baselineskip}

\begin{spacing}{1.0}
{\small\noindent\textbf{Note to the reader.}\enspace Problems from Walter Rudin are taken from
\textit{Principles of Mathematical Analysis}, 3rd~ed.\ (McGraw-Hill, 1976); the remaining problems
come from course homework, practice material, and exams. A {\color{green!60!black}\textbf{green}}
left rule marks an \textbf{assigned} problem, and a {\color{hwUnassignedBlue}\textbf{blue}} left rule
marks an \textbf{unassigned} problem included for completeness. Each solution is followed by a gray
\textsc{Reflections} block; these reflections are the student's study notes and are \emph{not} part
of the official submitted solution.\par}
\end{spacing}
\vspace{0.4\baselineskip}

\end{document}
